% Complete rate proofs; no characteristic-function or central-limit premise.
\section{Strong-law rates and normalization: proofs}
\label{sec:annex-p15-strong-law-rates-and-normalization}

The statements are in Chapter~\ref{sec:p15-strong-law-rates-and-normalization}.
The selected source proofs are Durrett's Theorems 2.5.11--2.5.13,
with the moment converse of Exercise 2.5.2
\cite[pp.~87--90]{DurrettPTEAuthor2019}.

\begin{annexitem}{annex:p15.power-tail-tools}{thm:p15-power-tail-tools}{Power integration, integer tails and decreasing sums}
\begin{proof}
For \(0\leq w\leq\infty\) and \(p>0\), the elementary improper integral
\(\int_0^w p y^{p-1}\,dy=w^p\) includes \(w=0\); at \(w=\infty\)
the equality follows by increasing the finite upper endpoint.
Consequently, pointwise,
\[
 W^p=\int_0^\infty p y^{p-1}\mathbf1_{\{y<W\}}\,dy.
\]
Assign the integrand value zero at \(y=0\). Away from that singleton
it is a nonnegative measurable function on the product of the
probability space and the Borel Lebesgue half-line. Its subgraph
indicator is measurable as in Theorem~\ref{thm:p15-layer-cake}.
Tonelli therefore gives the stated power-moment identity, even
when the expectation is infinite and even when \(p<1\).
For a finite nonnegative real \(v\),
\[
 \sum_{n\geq1}\mathbf1_{\{v>n\}}
       \leq v\leq1+\sum_{n\geq1}\mathbf1_{\{v>n\}}.
\]
This follows by counting the positive integers strictly below \(v\);
the same inequalities include integer \(v\) and \(v=0\).
At \(v=\infty\), both the count and \(v\) are infinite, so the
inequalities remain valid in the extended order.
Apply them to \(v=W^p\) and integrate the nonnegative sums to
obtain both tail-series bounds.

For \(q>1\), the function \(x^{-q}\) is positive decreasing.
Thus
\[
 \sum_{m=N}^\infty m^{-q}
  \leq N^{-q}+\int_N^\infty x^{-q}\,dx
  =N^{-q}+\frac{N^{1-q}}{q-1}
  \leq\frac{q}{q-1}N^{1-q}.
\]
If \(N\geq2\), comparison of each summand with the integral over
\([m-1,m]\) gives the other bound
\(\int_{N-1}^\infty x^{-q}\,dx=(N-1)^{1-q}/(q-1)\).
Likewise, for \(p>1\), the decreasing function \(x^{1/p-1}\)
gives
\[
 \sum_{m=1}^n m^{1/p-1}
 \leq1+\int_1^n x^{1/p-1}\,dx
 =1+p(n^{1/p}-1)\leq p n^{1/p}.
\]
Finally \(|u+v|\leq2\max(|u|,|v|)\), so raising to a positive
power gives the stated bound by \(2^p(|u|^p+|v|^p)\).
Apply it in both directions, writing \(X=(X-c)+c\) and
\(X-c=X+(-c)\), to prove the constant-shift moment assertion.
\end{proof}
\end{annexitem}

\begin{annexitem}{annex:p15.finite-variance-log-rate}{thm:p15-finite-variance-log-rate}{The exact variance-series argument}
\begin{proof}
Fix \(\varepsilon>0\), and set
\(a_n=\sqrt n(\log n)^{1/2+\varepsilon}\) for \(n\geq2\).
Choose \(a_1=a_2>0\). The sequence is positive, nondecreasing
and tends to infinity, since both positive factors increase from
\(n=2\) onward. The scaled variables \(X_n/a_n\) are independent,
centered and have variance \(\sigma^2/a_n^2\).
For \(x\geq2\), \(1/[x(\log x)^{1+2\varepsilon}]\) decreases,
and
\[
 \int_2^\infty\frac{dx}{x(\log x)^{1+2\varepsilon}}
       =\frac{(\log2)^{-2\varepsilon}}{2\varepsilon}<\infty.
\]
This follows by differentiating
\(-(\log x)^{-2\varepsilon}/(2\varepsilon)\).
Integral comparison, with the first summand treated separately,
therefore proves
\[
 \sum_{n\geq1}\operatorname{Var}(X_n/a_n)
  =\sigma^2\left(a_1^{-2}
       +\sum_{n=2}^\infty\frac1{n(\log n)^{1+2\varepsilon}}\right)
  <\infty.
\]
Theorem~\ref{thm:p15-variance-series-convergence} gives almost-sure
convergence of \(\sum_nX_n/a_n\). Apply
Theorem~\ref{thm:p15-kronecker-lemma} pointwise on that convergence
event to obtain \(S_n/a_n\to0\).
The argument includes \(\sigma^2=0\); then the displayed variance
sum is zero. For an i.i.d. sequence with finite second moment and
mean \(\mu\), apply the centered statement to \(X_n-\mu\).
Its finite second moment follows by the elementary square bound,
and its partial sum is \(S_n-n\mu\).
\end{proof}
\end{annexitem}

\begin{annexitem}{annex:p15.mz-source-rate}{thm:p15-mz-source-rate}{The complete source truncation, variance and mean proof}
\begin{proof}
Write \(W=|X_1|\), and use the power cutoff
\(Y_m=X_m\mathbf1_{\{|X_m|\leq m^{1/p}\}}\).
Identical laws and the integer-tail bound give
\[
 \sum_{m\geq1}\mathbb P(X_m\ne Y_m)
   =\sum_{m\geq1}\mathbb P(W^p>m)
   \leq\mathbb EW^p<\infty.
\]
The first Borel--Cantelli lemma implies only finitely many changed
summands almost surely. Since \(n^{1/p}\to\infty\), the difference
between the original and truncated partial sums divided by \(n^{1/p}\)
tends to zero on that event.

Here is the full variance calculation in the source's interval order.
For the nonnegative integer \(j-1\), let
\(I_j=[(j-1)^{1/p},j^{1/p})\), \(j\geq1\).
The second-moment tail integral and the bound
\(\mathbb P(|Y_m|>y)\leq\mathbf1_{\{y<m^{1/p}\}}\mathbb P(W>y)\)
give
\[
 \begin{aligned}
 \sum_{m\geq1}\frac{\operatorname{Var}(Y_m)}{m^{2/p}}
 &\leq\sum_{m\geq1}\frac{\mathbb EY_m^2}{m^{2/p}}\\
 &\leq\sum_{m\geq1}\sum_{j=1}^m
       \int_{I_j}\frac{2y}{m^{2/p}}\mathbb P(W>y)\,dy\\
 &=\sum_{j\geq1}\int_{I_j}
       2y\mathbb P(W>y)\sum_{m=j}^\infty m^{-2/p}\,dy .
 \end{aligned}
\]
All integrands are nonnegative and measurable. Tonelli for the
counting and Borel Lebesgue measures justifies each interchange;
changing interval endpoints in the Lebesgue integrals changes nothing.

Put \(q=2/p>1\). For \(j\geq2\), the source's integral comparison is
\[
 \sum_{m=j}^\infty m^{-2/p}
   \leq\frac{p}{2-p}(j-1)^{(p-2)/p}.
\]
To include the first interval with an explicit single constant, also
use the bound
\(\sum_{m=j}^\infty m^{-q}\leq[q/(q-1)]j^{1-q}\).
For \(y\in I_j\), \(j\geq\max(1,y^p)\), so this is at most
\([2/(2-p)]\max(1,y^p)^{1-2/p}\).
If \(0<y<1\), then \(y\leq y^{p-1}\), since \(p<2\).
If \(y\geq1\), multiplying the last bound by \(2y\) gives
\([4/(2-p)]y^{p-1}\) directly. In both ranges,
\[
 2y\sum_{m=j}^\infty m^{-2/p}
                   \leq\frac4{2-p}y^{p-1}\quad(y\in I_j,\ y>0).
\]
Insert this estimate in the integral and apply the power-tail identity:
\[
 \sum_{m\geq1}\frac{\mathbb EY_m^2}{m^{2/p}}
   \leq\frac4{2-p}\int_0^\infty y^{p-1}\mathbb P(W>y)\,dy
   =\frac4{p(2-p)}\mathbb EW^p<\infty.
\]

The variables \((Y_m-\mu_m)/m^{1/p}\) are mutually independent,
centered, and have summable variances. Independent-series convergence
and Kronecker's lemma with weights \(m^{1/p}\) therefore imply
\[
 n^{-1/p}\sum_{m=1}^n(Y_m-\mu_m)\longrightarrow0
                                      \quad\text{almost surely}.
\]
It remains to estimate the means; the centered-series conclusion
does not by itself give the uncentered rate.
Since \(\mathbb EX_1=0\), integration over the common law gives
\(\mu_m=-\mathbb E[X_1\mathbf1_{\{W>m^{1/p}\}}]\).
On \(\{W>m^{1/p}\}\), \(W/m^{1/p}>1\), so
\[
 \begin{aligned}
 |\mu_m|
 &\leq m^{1/p}\mathbb E\!\left[
       (W/m^{1/p})\mathbf1_{\{W>m^{1/p}\}}\right]\\
 &\leq m^{1/p}\mathbb E\!\left[
       (W/m^{1/p})^p\mathbf1_{\{W>m^{1/p}\}}\right]\\
 &=m^{1/p-1}d_m,\qquad
 d_m=\mathbb E[W^p\mathbf1_{\{W>m^{1/p}\}}].
 \end{aligned}
\]
Dominated convergence, using the integrable \(W^p\), gives \(d_m\to0\).
Given \(\delta>0\), choose \(M\) so that \(d_m\leq\delta\)
for every \(m\geq M\). The numerical sum bound then yields
\[
 n^{-1/p}\sum_{m=1}^n|\mu_m|
  \leq n^{-1/p}\sum_{m=1}^{M-1}|\mu_m|
       +\delta n^{-1/p}\sum_{m=M}^n m^{1/p-1}
  \leq n^{-1/p}\sum_{m=1}^{M-1}|\mu_m|+p\delta .
\]
The finite head tends to zero; letting \(\delta\downarrow0\)
proves the required mean rate. Adding it to the centered rate and
then restoring the finitely many changed summands proves the theorem.
\end{proof}
\end{annexitem}

\begin{annexitem}{annex:p15.mz-subunit-rate}{thm:p15-mz-subunit-rate}{Absolute weighted summability below order one}
\begin{proof}
Set \(W=|X_1|\) and \(Y_n=X_n\mathbf1_{\{|X_n|\leq n^{1/p}\}}\).
The same tail-probability estimate gives
\(\sum_n\mathbb P(X_n\ne Y_n)\leq\mathbb EW^p<\infty\).
It uses only the common law and the first Borel--Cantelli lemma.
For the nonnegative weighted expectations, Tonelli and the common
law give
\[
 \sum_{n\geq1}\frac{\mathbb E|Y_n|}{n^{1/p}}
    =\mathbb E\!\left[
       W\sum_{\substack{n\geq1\\ n\geq W^p}}n^{-1/p}\right].
\]
For a fixed \(w\geq0\), put \(N=\max(1,\lceil w^p\rceil)\)
and \(q=1/p>1\). The numerical bound in
Lemma~\ref{thm:p15-power-tail-tools} bounds the inner sum by
\([1/(1-p)]N^{1-1/p}\).
Since \(N\geq\max(1,w^p)\), multiplication by \(w\) makes this
at most \([1/(1-p)]w^p\): for \(w\geq1\), the exponents give
\(w(w^p)^{1-1/p}=w^p\); for \(0\leq w<1\), use \(w\leq w^p\).
Thus
\[
 \mathbb E\!\left[\sum_{n\geq1}|Y_n|/n^{1/p}\right]
                  \leq\frac{\mathbb E|X_1|^p}{1-p}<\infty.
\]
A nonnegative measurable function with finite integral is finite
almost surely. Hence this weighted absolute series converges almost
surely. The first Borel--Cantelli conclusion changes only finitely
many finite weighted terms, so \(\sum_n|X_n|/n^{1/p}<\infty\)
almost surely as well. On that common event, apply Kronecker's lemma
to the numerical summands \(|X_n(\omega)|\) and weights \(n^{1/p}\).
It yields the asserted normalized absolute-sum limit. The inequality
\(|S_n|\leq\sum_{m=1}^n|X_m|\) gives the signed limit. No
independence was used anywhere in this argument.
\end{proof}
\end{annexitem}

\begin{annexitem}{annex:p15.power-rate-necessity}{thm:p15-power-rate-necessity}{Normalized increments force the moment, and averages force the center}
\begin{proof}
Put \(V_n=X_n-c\) and \(\Delta_n=\sum_{j=1}^nV_j=S_n^{(c)}\).
At any outcome where \(\Delta_n/n^{1/p}\to0\), for \(n\geq2\),
\[
 \frac{V_n}{n^{1/p}}
  =\frac{\Delta_n}{n^{1/p}}
       -\left(\frac{n-1}{n}\right)^{1/p}
                 \frac{\Delta_{n-1}}{(n-1)^{1/p}}
                         \longrightarrow0.
\]
If \(\mathbb E|V_1|^p=\infty\), the upper integer-tail estimate
forces \(\sum_n\mathbb P(|V_1|^p>n)=\infty\).
The events \(\{|V_n|>n^{1/p}\}\) are independent and have exactly
these probabilities. The second Borel--Cantelli lemma makes them
occur infinitely often almost surely, contradicting the necessary
increment limit. Thus the zero-rate event has probability zero when
this moment is infinite. Constant-shift moment equivalence makes
\(\mathbb E|V_1|^p<\infty\) equivalent to
\(\mathbb E|X_1|^p<\infty\), giving both assertions about the moment.

If \(p\geq1\), the inequality \(|x|\leq1+|x|^p\) implies
integrability once that moment is finite. The integrable strong law
therefore gives \(S_n/n\to\mathbb EX_1\) almost surely.
On the assumed zero-rate event,
\[
 S_n/n-c=n^{1/p-1}\bigl(S_n^{(c)}/n^{1/p}\bigr)\longrightarrow0,
\]
since \(0<n^{1/p-1}\leq1\). On the intersection of the two
probability-one events the deterministic limiting values must agree,
so \(c=\mathbb EX_1\).
\end{proof}
\end{annexitem}

\begin{annexitem}{annex:p15.mz-full-classification}{thm:p15-mz-full-classification}{Combining the three parameter regimes}
\begin{proof}
Necessity in both regimes follows from the preceding lemma.
For \(0<p<1\), the subunit theorem gives \(S_n/n^{1/p}\to0\)
under the moment condition. For any fixed finite \(c\),
\(nc/n^{1/p}=c n^{1-1/p}\to0\); subtracting it gives the same
rate with that center.
At \(p=1\), the moment condition is \(\mathbb E|X_1|<\infty\),
and the integrable strong law gives
\((S_n-n\mathbb EX_1)/n\to0\).
For \(1<p<2\), the moment condition gives a finite mean \(\mu\).
The shift bound shows that \(X_n-\mu\) still has finite \(p\)-th
absolute moment; it is centered and i.i.d. Apply the full source
theorem to this sequence to obtain
\((S_n-n\mu)/n^{1/p}\to0\).
These sufficiency arguments and the already proved necessity give
exactly the two stated equivalences, without including the endpoint
two or assigning a nonexistent mean in the subunit range.
\end{proof}
\end{annexitem}

\begin{annexitem}{annex:p15.feller-infinite-mean}{thm:p15-feller-infinite-mean}{Both full Feller arguments and every truncation limit}
\begin{proof}
Write \(b_n=a_n/n>0\). Its monotonicity implies
\(a_{n+1}\geq((n+1)/n)a_n>a_n\), \(a_n\geq n a_1\),
and \(a_{kn}\geq k a_n\) for positive integers \(k,n\).
Thus the \(a_n\) are strictly increasing and diverge.
Also \(q_n=\mathbb P(W\geq a_n)\) is nonincreasing.

Suppose first that \(Q=\infty\). For a fixed positive integer \(k\),
\[
 \begin{aligned}
 \sum_{n\geq1}\mathbb P(W\geq k a_n)
 &\geq\sum_{n\geq1}\mathbb P(W\geq a_{kn})\\
 &\geq\frac1k\sum_{m=k}^\infty q_m=\infty.
 \end{aligned}
\]
For the second inequality, partition \(m\geq k\) into blocks
\(kn\leq m<k(n+1)\); each of their \(k\) terms \(q_m\) is at
most \(q_{kn}\). Independence and the second Borel--Cantelli lemma
give \(|X_n|\geq k a_n\) infinitely often almost surely.
Intersect these events over integers \(k\geq1\). On that common
event \(\limsup_n|X_n|/a_n=\infty\).
Since \(X_n=S_n-S_{n-1}\),
\(\max(|S_n|,|S_{n-1}|)\geq|X_n|/2\).
If \(n\geq2\) and the large-jump event at level \(k\) occurs,
then either \(|S_n|/a_n\geq k/2\), or
\[
 |S_{n-1}|/a_{n-1}\geq|S_{n-1}|/a_n\geq k/2.
\]
Infinitely many jump indices give infinitely many indices in these
two possible positions; a finite set of sum indices could account
for only finitely many jumps. Thus the sum limit superior is at
least \(k/2\) for every \(k\), proving the divergent branch.

Now suppose \(Q<\infty\). Let \(\nu\) be the Borel probability
law of \(W=|X_1|\) on \([0,\infty)\), and set
\(A_m=[a_{m-1},a_m)\), \(m\geq1\), with \(a_0=0\).
These intervals partition the nonnegative half-line. If a finite
\(w\) belongs to \(A_m\), exactly the first \(m\) thresholds
\(a_0,\ldots,a_{m-1}\) are at most \(w\). Therefore pointwise
\[
 \sum_{m\geq1}m\mathbf1_{A_m}(w)
       =\sum_{n\geq1}\mathbf1_{\{w\geq a_{n-1}\}}.
\]
Integrating the nonnegative sums gives the exact source count identity
\[
 \sum_{m\geq1}m\nu(A_m)=
       \sum_{n\geq1}\mathbb P(W\geq a_{n-1})=1+Q<\infty.
\]
The initial one includes possible mass at \(W=0\).

The strict truncations \(Y_n=X_n\mathbf1_{\{|X_n|<a_n\}}\)
differ from \(X_n\) with probability \(q_n\). Thus only finitely
many changes occur almost surely by the first Borel--Cantelli lemma.
Each \(Y_n\) has finite second moment. The full source variance
calculation, using the law of the absolute value, is
\[
 \begin{aligned}
 \sum_{n\geq1}\operatorname{Var}(Y_n/a_n)
 &\leq\sum_{n\geq1}a_n^{-2}
                   \sum_{m=1}^n\int_{A_m}y^2\,d\nu(y)\\
 &=\sum_{m\geq1}\int_{A_m}y^2\,d\nu(y)
                              \sum_{n=m}^\infty a_n^{-2}.
 \end{aligned}
\]
Tonelli justifies the nonnegative interchange. For \(n\geq m\),
\(a_n\geq n a_m/m\), and the reciprocal-square bound gives
\[
 \sum_{n=m}^\infty a_n^{-2}
   \leq\frac{m^2}{a_m^2}\sum_{n=m}^\infty n^{-2}
   \leq\frac{2m}{a_m^2}.
\]
On \(A_m\), \(y<a_m\), so
\[
 \sum_{n\geq1}\operatorname{Var}(Y_n/a_n)
          \leq2\sum_{m\geq1}m\nu(A_m)=2(1+Q)<\infty.
\]
Independent-series convergence for \((Y_n-\mathbb EY_n)/a_n\),
followed by Kronecker's lemma, now gives
\[
 \frac{T_n-\mathbb ET_n}{a_n}\longrightarrow0
                                      \quad\text{almost surely}.
\]

To complete the source's mean argument, first prove \(a_n/n\to\infty\).
If its nondecreasing values were bounded above by a finite \(B>0\),
then \(a_n\leq Bn\) and
\[
 \sum_{n\geq1}\mathbb P(W/B>n)
   \leq\sum_{n\geq1}\mathbb P(W\geq Bn)
   \leq Q<\infty.
\]
The integer-tail upper bound would give \(\mathbb E(W/B)<\infty\),
contradicting \(\mathbb EW=\infty\). Thus the ratio diverges.

For a fixed integer \(N\geq1\) and \(n>N\), finite signed
linearity, common laws and increasing cutoffs give
\[
 \begin{aligned}
 \left|\frac1{a_n}\sum_{m=1}^n\mathbb EY_m\right|
 &\leq\frac1{a_n}\sum_{m=1}^n
                  \mathbb E[W\mathbf1_{\{W<a_m\}}]\\
 &\leq\frac{n}{a_n}\mathbb E[W\mathbf1_{\{W<a_n\}}]\\
 &\leq\frac{n a_N}{a_n}
       +\frac{n}{a_n}\mathbb E[W\mathbf1_{\{a_N\leq W<a_n\}}].
 \end{aligned}
\]
Every truncated integral here is finite. The first term tends to
zero for fixed \(N\), because \(a_n/n\to\infty\).
For the second term use the same intervals, and
\(n/a_n\leq m/a_m\) whenever \(m\leq n\):
\[
 \begin{aligned}
 \frac{n}{a_n}\mathbb E[W\mathbf1_{\{a_N\leq W<a_n\}}]
 &=\sum_{m=N+1}^n\frac{n}{a_n}\int_{A_m}y\,d\nu(y)\\
 &\leq\sum_{m=N+1}^n\frac{m}{a_m}\int_{A_m}y\,d\nu(y)\\
 &\leq\sum_{m=N+1}^\infty m\nu(A_m).
 \end{aligned}
\]
This is the tail of the finite series in the count identity, so it
tends to zero as \(N\to\infty\). First taking the limit superior
in \(n\), then \(N\to\infty\), proves \(\mathbb ET_n/a_n\to0\).
Adding the centered limit gives \(T_n/a_n\to0\) almost surely.
Restoring the finitely many changed summands changes \(T_n\) by
an eventually constant finite outcome-dependent number; division
by \(a_n\to\infty\) removes it. Hence \(S_n/a_n\to0\), completing
the finite-tail branch and the whole dichotomy.
\end{proof}
\end{annexitem}

\begin{annexitem}{annex:p15.rate-calculated-examples}{thm:p15-rate-calculated-examples}{Computing the moments and the two threshold sums}
\begin{proof}
Fair signs are centered and have second moment one. Taking
\(\varepsilon=1/2\) in the finite-variance theorem gives the first
displayed normalization \(\sqrt n\log n\).

For the first positive dyadic law, put \(u=2^{1/4}>1\).
Its masses are \((u^3-1)u^{-3j}\), whose sum is one by the
geometric-series formula. Its square-root moment is
\[
 (u^3-1)\sum_{j\geq1}u^{-j}
        =\frac{u^3-1}{u-1}=u^2+u+1
        =\sqrt2+2^{1/4}+1.
\]
Its first-moment series is
\((u^3-1)\sum_{j\geq1}u^j=\infty\).
The subunit theorem at \(p=1/2\) therefore gives \(S_n/n^2\to0\);
the positive infinite-mean strong law gives \(S_n/n\to\infty\).
Intersect the two probability-one events to obtain both limits together.

For the proposed density, direct integration gives
\[
 \int_1^\infty\frac74 x^{-11/4}\,dx=1,\qquad
 \int_1^\infty\frac74 x^{-7/4}\,dx=\frac73,\qquad
 \int_1^\infty\frac74 x^{-5/4}\,dx=7.
\]
These are respectively the normalization, first moment and
\(3/2\)-moment. For the second moment the integrand is
\((7/4)x^{-3/4}\), whose integral diverges. All these integrals
are nonnegative limits of elementary bounded-interval integrals.
The constant-shift inequality proves finiteness of
\(\mathbb E|X_1-7/3|^{3/2}\). Apply the centered source
Marcinkiewicz--Zygmund theorem at \(p=3/2\) to obtain the stated
\(n^{2/3}\) rate.

For the last dyadic law the masses sum to one and its mean is
\(\sum_{j\geq1}2^j2^{-j}=\infty\).
For every real \(x>0\), let \(j_0\geq1\) be the least integer
with \(2^{j_0}\geq x\). Its inclusive tail satisfies
\[
 \mathbb P(X_1\geq x)=\sum_{j\geq j_0}2^{-j}
                     =2^{1-j_0}\leq2/x.
\]
Thus \(\sum_n\mathbb P(X_1\geq n^2)\leq2\sum_nn^{-2}<\infty\).
For \(a_n=n\), the tail series is infinite, either by the
integer-tail moment criterion or by summing dyadic blocks.
Both \(a_n/n=1\) and \(a_n/n=n\) are nondecreasing. Feller's
two branches give the two claimed limits. The specified i.i.d.
sequences exist by the already proved countable independent-sequence
construction.
\end{proof}
\end{annexitem}

\begin{annexitem}{annex:p15.rate-hypothesis-boundaries}{thm:p15-rate-hypothesis-boundaries}{The order-two bound and each explicit hypothesis failure}
\begin{proof}
For independent fair signs, apply the bounded fourth-moment lemma
with bound \(M=1\), sum \(S_n\) and variance \(v=n\).
It gives
\[
 \mathbb P(|S_n|>\sqrt{n/2})
      \geq\frac{n}{4(1+3n)}\geq\frac1{16}.
\]
The last inequality is equivalent to \(n\geq1\).
This is exactly the asserted probability bound, which rules out
convergence in probability of \(S_n/\sqrt n\) to zero.
Let \(C=\{S_n/\sqrt n\to0\}\). It is measurable by the
countable threshold criterion for convergence to zero.
Removing any fixed initial block changes the normalized sums by
a fixed finite value divided by \(\sqrt n\), which tends to
zero. Expressing the same convergence event with tail sums shows
that \(C\) belongs to every tail-coordinate sigma-algebra.
The independent tail zero-one law therefore gives
\(\mathbb P(C)\in\{0,1\}\).
If its probability were one, the indicators of the deviation
events just bounded would tend to zero almost surely; bounded
convergence would make their probabilities tend to zero, a
contradiction. Thus \(\mathbb P(C)=0\).
For \(p>2\), a zero \(n^{1/p}\) rate would imply a zero
square-root rate after multiplication by \(n^{1/p-1/2}\).
Its event is consequently contained in \(C\) and is null as well.

The deterministic sequence \(X_n=1\) is independent and has every
absolute moment finite. Its power-normalized sum is
\(n^{1-1/p}\), which equals one at \(p=1\) and diverges for
\(p>1\). To check the logarithmic claim too, put
\(\beta=1/2+\varepsilon>0\) and choose \(d=1/(4\beta)>0\).
For \(n\geq2\),
\[
 \log n=\int_1^n t^{-1}\,dt
       \leq\int_1^n t^{d-1}\,dt
       \leq n^d/d.
\]
It follows that
\(\sqrt n/(\log n)^\beta\geq d^\beta n^{1/4}\to\infty\).
Thus centering is a substantive condition even when the variance
of each summand is zero.

For the critical-moment dyadic law,
\[
 \sum_{j\geq1}(2^p-1)2^{-pj}=1,\qquad
 \mathbb E X_1^p=(2^p-1)\sum_{j\geq1}1=\infty.
\]
The moment-necessity lemma gives probability zero for a zero rate
with any finite center. When \(p>1\), its first moment is the
finite geometric sum
\[
 \mathbb EX_1=\frac{2^p-1}{2^{p-1}-1}.
\]
Thus even its proper finite center cannot compensate for the missing
\(p\)-th moment.

For the finite-mean Feller counterexample, with \(X_n=1\) and
\(a_n=n\), the inclusive threshold probability is one at \(n=1\)
and zero afterward. Hence \(Q=1\), while the normalized sum is
one. All the normalization and independence assumptions hold;
the infinite-first-absolute-moment assumption has been removed.
For the shared dyadic variable, its law has infinite first moment
and the threshold series at \(a_n=n\) diverges. But
\(S_n=nX\), so every normalized sum equals the finite value
\(X(\omega)\). These variables are not independent: for example
the events that \(X_1\) and \(X_2\) equal two coincide and have
probability \(1/2\), not joint probability \(1/4\).
This verifies the stated dependence failure.
\end{proof}
\end{annexitem}

\begin{annexitem}{annex:p15.feller-ratio-counterexample}{thm:p15-feller-ratio-counterexample}{A deterministic staircase with summable individual tails}
\begin{proof}
Use i.i.d. positive variables with
\(\mathbb P(X_n=2^j)=2^{-j}\), \(j\geq1\).
Their mean is infinite, and the extended-mean strong law already
proved gives \(S_n/n\to\infty\) almost surely.
For each fixed integer \(k\geq1\), bounded convergence of
\(\mathbf1_{\{S_n/n\leq k^3\}}\) therefore gives
\(\mathbb P(S_n/n\leq k^3)\to0\).
Consequently the set of integers required in the recursive definition
of \(N_k\) is nonempty, even after requiring \(n\geq N_{k-1}+2\).
It has a least member. Every \(N_k\) is thus a finite deterministic
integer determined by the law and \(k\), not by a realized sample.
They increase by at least two and tend to infinity.

Put \(A_k=N_k k^2\) and \(a_n=A_k\) on the indicated block.
The values \(A_k\) strictly increase, so \(a_n\) is positive
and nondecreasing and tends to infinity. Each block contains at
least two consecutive indices with the same positive denominator;
the ratios \(a_n/n\) strictly decrease at those steps.
Thus the exact Feller ratio condition fails, even though the
denominators themselves are nondecreasing.

The inclusive dyadic tail bound proved in the preceding calculations
gives \(\mathbb P(X_1\geq x)\leq2/x\) for all \(x>0\).
Therefore
\[
 \sum_{n\geq1}\mathbb P(X_1\geq a_n)
  =\sum_{k\geq1}(N_k-N_{k-1})\mathbb P(X_1\geq A_k)
  \leq\sum_{k\geq1}\frac{2N_k}{N_k k^2}
  =2\sum_{k\geq1}k^{-2}<\infty.
\]
By construction,
\(\mathbb P(S_{N_k}/N_k\leq k^3)\leq2^{-k}\).
These error probabilities are summable, so the first Borel--Cantelli
lemma says that almost surely all sufficiently large \(k\) satisfy
\(S_{N_k}/N_k>k^3\). On that event,
\[
 \frac{S_{N_k}}{a_{N_k}}
       =\frac{S_{N_k}/N_k}{k^2}>k
 \qquad\text{for every sufficiently large }k.
\]
The limit superior along this subsequence, and hence along the full
sequence, is infinite. This proves every claimed property of the
counterexample using only the already established strong law and
Borel--Cantelli; no distributional limit theorem is assumed.
\end{proof}
\end{annexitem}
